\section{Descrizione e contesto}

\subsection{Contesto}
Un'azienda intende sviluppare un sistema per raccogliere e consultare \emph{testi
musicali (lyrics)} associati a brani di qualsiasi artista. L'obiettivo è mantenere
un database di testi in formato sia piano (\emph{plain lyrics}) sia sincronizzato
con i \emph{timestamp} della traccia (formato \emph{LRC}), consentendo agli utenti
di cercare i testi tramite titolo, artista e album, e di contribuire
pubblicandone di nuovi.

Per limitare pubblicazioni automatiche abusive da parte di utenti non
autenticati, il sistema richiede la risoluzione di una sfida crittografica
(\emph{Proof of Work}, PoW) prima di ogni pubblicazione anonima. Gli utenti
registrati possono invece pubblicare senza PoW, con il brano attribuito al
loro account.

\subsection{Obiettivo del documento}
Il presente documento descrive in modo formale i requisiti del sistema Stopify:
stakeholder coinvolti, requisiti funzionali e non funzionali, architettura
hardware/software e di rete, casi d'uso e analisi degli aspetti di sicurezza.
Costituisce la base per lo sviluppo, il collaudo e la futura evoluzione del
prodotto.

\subsection{Assunzioni}
\begin{itemize}
    \item Il backend è una \emph{Web API REST} sviluppata in Python con il
    framework Flask. La persistenza è demandata a un database SQLite locale.
    \item Il sistema espone gli endpoint richiesti dalla specifica
    \emph{LRCLIB} (\url{https://lrclib.net/docs}), opportunamente emulati
    sulla cache locale.
    \item È sviluppata anche un'applicazione mobile cross-platform (Android oriented) basata su React Native ed Expo che funge da client
    principale per la consultazione e la pubblicazione.
    \item L'autenticazione è basata su \emph{JSON Web Token} (JWT):
    il server emette un \emph{access token} al login/registrazione,
    che il client allega ad ogni richiesta protetta.
    \item L'ambiente di esecuzione è una rete locale fidata; tutte le
    considerazioni che riguardano deployment pubblico (TLS, rate limiting,
    backup, monitoraggio) sono trattate come \emph{future work}.
\end{itemize}
